%% Question bank chapter file - generated by bank_tools.py build (do not edit while the bank is still being built:
%% it is regenerated from the per-test data in Question Bank\data). Every question is one bankq block:
%%   \begin{bankq}[id=..., source={...}, paper=.., qno=.., type=..., subtopic={...}, answer={...}, answer_src=key|solved,
%%                 working={...}, difficulty=1-10, diff_parts={K S I C F}, diff_note={...},
%%                 group=<passage id or empty>, also_in={duplicate copies in other tests}]
%% difficulty = min(10, K+S+I+C+F-1): K concepts 1-3, S steps 1-4, I insight 0-2, C computation 0-2, F format 0-1
%%   ...question LaTeX, options, figures (\bfig / \bfigi -> img/)...
%%   \end{bankq}
%% Passages shared by several questions are bankpassage blocks; their questions carry group=<passage id>.
\documentclass[10pt]{article}
\usepackage[a4paper,margin=1.8cm]{geometry}
\usepackage[T1]{fontenc}
\usepackage{amsmath,amssymb,mathtools}
\usepackage{graphicx}
\usepackage[export]{adjustbox}
\usepackage{xcolor,needspace,keyval}
\usepackage{enumitem}
\usepackage{booktabs,array,tabularx}
\usepackage[hidelinks]{hyperref}
\setlength{\parindent}{0pt}
\setlength{\parskip}{3pt}
\makeatletter
\define@key{bq}{id}{\def\bq@id{#1}}\define@key{bq}{source}{\def\bq@src{#1}}\define@key{bq}{paper}{\def\bq@paper{#1}}
\define@key{bq}{qno}{\def\bq@qno{#1}}\define@key{bq}{type}{\def\bq@type{#1}}\define@key{bq}{subtopic}{\def\bq@sub{#1}}
\define@key{bq}{answer}{\def\bq@ans{#1}}\define@key{bq}{answer_src}{\def\bq@asrc{#1}}\define@key{bq}{working}{\def\bq@work{#1}}
\define@key{bq}{group}{\def\bq@group{#1}}\define@key{bq}{also_in}{\def\bq@also{#1}}
\define@key{bq}{difficulty}{\def\bq@diff{#1}}\define@key{bq}{diff_parts}{\def\bq@dparts{#1}}\define@key{bq}{diff_note}{\def\bq@dnote{#1}}
\newenvironment{bankq}[1][]{\def\bq@work{}\def\bq@also{}\def\bq@group{}\def\bq@diff{?}\def\bq@dparts{}\def\bq@dnote{}\setkeys{bq}{#1}%
  \par\needspace{4\baselineskip}\medskip\noindent{\small\textbf{\bq@id}\hfill\color{gray}[\bq@src, Paper \bq@paper, Q\bq@qno]
  \bq@type\ $\cdot$ \bq@sub\ $\cdot$ \textbf{D\bq@diff}}\par\nopagebreak\relax}
  {\par{\small\color{teal}Answer: \textbf{\bq@ans}\ifx\bq@work\@empty\else\ (solved: \bq@work)\fi
  \ifx\bq@also\@empty\else\par\color{gray}Also in: \bq@also\fi}\par}
\define@key{bp}{id}{\def\bp@id{#1}}\define@key{bp}{source}{\def\bp@src{#1}}\define@key{bp}{questions}{\def\bp@qs{#1}}
\newenvironment{bankpassage}[1][]{\setkeys{bp}{#1}\par\needspace{4\baselineskip}\bigskip\noindent
  {\small\textbf{Passage \bp@id}\hfill\color{gray}[\bp@src; questions \bp@qs]}\par\nopagebreak\relax\itshape}{\par}
\makeatother
\newcommand{\bfig}[2]{\par\smallskip{\centering\includegraphics[width=#1pt,max width=0.62\linewidth,max height=0.32\textheight]{img/#2}\par}\smallskip}
\newcommand{\bfigi}[2]{\includegraphics[width=#1pt,max width=\linewidth,max height=0.32\textheight,valign=c]{img/#2}}
\newcommand{\bmissing}[1]{\textit{[figure missing in the source conversion: \detokenize{#1}]}}
\newenvironment{qoptions}{\begin{itemize}[leftmargin=1.8em,itemsep=1pt,topsep=2pt,label={}]}{\end{itemize}}
\newcommand{\opt}[2]{\item[(#1)] #2}
\providecommand{\fwstart}{}\providecommand{\fwend}{}\providecommand{\qfullsplit}{}
\newcolumntype{L}[1]{>{\raggedright\arraybackslash}p{#1}}
\newcolumntype{C}[1]{>{\centering\arraybackslash}p{#1}}
\newcolumntype{R}[1]{>{\raggedleft\arraybackslash}p{#1}}
\providecommand{\sectionhdr}[1]{}\providecommand{\subsectionhdr}[1]{}


\begin{document}
{\Large\bfseries COM --- Question Bank}\par\smallskip
{\small Physics chapter P04. Questions from all converted test papers, grouped by source series. Answers: printed key where it matches the paper, otherwise solved (answer\_src=solved, with working). Exact repeats across papers appear once, with also\_in listing the other copies.}\par\medskip\hrule

\section*{Score-I papers}
\begin{bankq}[id=SI12-P1-PHY-S3-Q08, source={Score-I - Test 12 (09 Mar 2022)}, paper=1, qno=8, type=numerical, subtopic={Collisions}, answer={4}, answer_src=key, difficulty=2, diff_parts={K1 S1 I0 C0 F1}, diff_note={concepts: elastic head-on collision | steps: momentum conservation with e = 1; solve for the speed of B}]
The velocities of two steel balls before impact are shown. After head on elastic collision find the velocity (in m/s) of ball $B$.
\bfig{130.8}{SI12_p1_phy_s3_q08_1.png}
\end{bankq}

\section*{TOAS papers}
\begin{bankq}[id=TOAS02-P1-PHY-S2-Q09, source={TOAS - Test 02 (13 Feb 2022)}, paper=1, qno=9, type=multi, subtopic={Collisions}, answer={B,C}, answer_src=key, difficulty=6, diff_parts={K3 S3 I0 C0 F1}, diff_note={concepts: elastic collision, SHM, projectile | steps: mass of block 2 from the period; velocities after the elastic collision; fall time and horizontal range of block 1; amplitude of block 2; check the options}]
Figure shows block 1 of mass 0.5 kg sliding to the right over a frictionless elevated surface at a speed of 12 m/s. The block undergoes an elastic collision with stationary block 2, which is attached to a spring of spring constant 100 N/m (Assume that spring doesn't affect collision). After the collision, block 2 oscillates in SHM with a time period of 0.628 s and block 1 slides to the opposite end of the elevated surface landing at a distance $d$ from the base of that surface after falling through a height $h=4.90$ m. Then which of the following is (are) \textbf{CORRECT}: (take $g=9.8\ \mathrm{m/s^2}$ and $\pi=3.14$)
\bfig{157.0}{TOAS02_p1_phy_s2_q09_1.png}
\begin{qoptions}
\opt{A}{The value of $d=2$ m}
\opt{B}{Spring potential energy is $<35$ J}
\opt{C}{Maximum kinetic energy of block 2 is 32 J}
\opt{D}{Speed of block 1 when it hits the ground is 4 m/s}
\end{qoptions}
\end{bankq}
\begin{bankq}[id=TOAS04-P1-PHY-S1-Q02, source={TOAS - Test 04 (15 Mar 2022)}, paper=1, qno=2, type=single, subtopic={Momentum}, answer={C}, answer_src=key, difficulty=2, diff_parts={K1 S2 I0 C0 F0}, diff_note={concepts: impulse and work-energy theorem | steps: impulse from the F-t graph gives velocities; check the velocity at t = 2 s; final velocity and work by the force}]
A particle of mass 2 kg is moving along $x$-axis with a velocity of 5 m/s. It crosses the origin at $t=0$. A time varying force whose variation is as shown starts acting on a particle at $t=1$ sec. For this situation, mark out the correct statement.
\bfig{118.8}{TOAS04_p1_phy_s1_q02_1.png}
\begin{qoptions}
\opt{A}{The particle will come to rest instantaneously at $t=2$ sec}
\opt{B}{The final velocity of the particle will be 18 m/s.}
\opt{C}{Work done by the force $F$ on the particle is 375 J.}
\opt{D}{All of the above}
\end{qoptions}
\end{bankq}
\begin{bankq}[id=TOAS07-P1-PHY-S1-Q01, source={TOAS - Test 07 (01 May 2022)}, paper=1, qno=1, type=single, subtopic={Variable mass}, answer={A}, answer_src=key, difficulty=2, diff_parts={K1 S1 I1 C0 F0}, diff_note={concepts: variable mass motion | steps: spilled sand leaves with the cart velocity so F = m dv/dt with m = m0 - mu t; integrate | I: spilled sand carries no relative momentum}]
A cart loaded with sand moves along a horizontal floor due to a constant force $F$ coinciding in direction with the cart's velocity vector. In the process sand spills through a hole. in the bottom with constant rate $\mu$ kg/sec. Find the velocity of the cart as a function of time `$t$', if at the initial moment $t=0$ the cart with loaded sand had the mass $m_0$ and its velocity was equal to zero.
\bfig{116.3}{TOAS07_p1_phy_s1_q01_1.png}
\begin{qoptions}
\opt{A}{$\dfrac{F}{\mu}\ln\!\left(\dfrac{m_0}{m_0-\mu t}\right)$}
\opt{B}{$\dfrac{F}{\mu}\ln\!\left(\dfrac{m_0}{m_0+\mu t}\right)$}
\opt{C}{$\dfrac{F}{\mu}\ln\!\left(\dfrac{m_0+\mu t}{m_0}\right)$}
\opt{D}{$\dfrac{F}{\mu}\ln\!\left(\dfrac{m_0-\mu t}{m_0}\right)$}
\end{qoptions}
\end{bankq}
\begin{bankq}[id=TOAS07-P1-PHY-S2-Q05, source={TOAS - Test 07 (01 May 2022)}, paper=1, qno=5, type=multi, subtopic={Centre of mass}, answer={B,C}, answer_src=key, difficulty=3, diff_parts={K1 S2 I0 C0 F1}, diff_note={concepts: centre of mass of a spring-block system | steps: centre of mass acceleration F/2m; displacement F t\^{}2/4m; positions of A and B from the extension}]
Two blocks $A$ and $B$ each of mass `$m$', are connected by a massless spring of force constant $k$. Initially given system is at rest and spring is in its natural length. A constant force $F$ starts acting on block $A$ at time $t=0$ and at time $t$, the extension in the spring is seen to be $\ell$. Given system lies on smooth horizontal surface.
\bfig{119.9}{TOAS07_p1_phy_s2_q05_1.png}
\begin{qoptions}
\opt{A}{Acceleration of centre of mass of the given system ($A+B+$ spring) is $\dfrac Fm$.}
\opt{B}{Displacement of centre of mass of the given system ($A+B+$ spring) in time $t$ is $\dfrac{Ft^2}{4m}$.}
\opt{C}{Displacement of block $A$ in time $t$ is $\dfrac{Ft^2}{4m}+\dfrac{\ell}{2}$.}
\opt{D}{Displacement of block $A$ in time $t$ is $\dfrac{Ft^2}{2m}-\dfrac{\ell}{2}$.}
\end{qoptions}
\end{bankq}
\begin{bankq}[id=TOAS07-P1-PHY-S2-Q10, source={TOAS - Test 07 (01 May 2022)}, paper=1, qno=10, type=multi, subtopic={Centre of mass}, answer={B,C,D}, answer_src=key, difficulty=5, diff_parts={K2 S2 I1 C0 F1}, diff_note={concepts: centre of mass frame, spring system | steps: velocity of the centre of mass; relative velocity and reduced mass; maximum extension from energy; time to first maximum compression | I: reduced mass for the two-body spring}]
Two blocks A(5kg) and B(2kg) attached to the ends of a spring of spring constant 1120 N/m are placed on a smooth horizontal plane with the spring undeformed. Simultaneously velocities of 3m/s and 10 m/s along the line of the spring in the same direction are imparted to A and B at $t=0$ then:
\bfig{119.5}{TOAS07_p1_phy_s2_q10_1.png}
\begin{qoptions}
\opt{A}{When the extension of the spring is maximum the velocities of A and B are zero.}
\opt{B}{The maximum extension of the spring is 25 cm.}
\opt{C}{The first maximum compression occurs at $t=\dfrac{3\pi}{56}$ second.}
\opt{D}{Maximum extension and maximum compression occur alternately.}
\end{qoptions}
\end{bankq}
\begin{bankpassage}[id=TOAS07-P2-PHY-S3-G3-4, source={TOAS - Test 07 (01 May 2022)}, questions={3-4}]
A system consists of two small spheres of mass 1kg and 2kg interconnected by a weightless spring. At $t=0$, the spheres are set in motion with initial velocities as shown in figure. Position vector of centre of mass at $t=0$ is $20\,\hat j$. Initially spring is in its natural length. System moves in earth's gravitational field.
\bfig{109.4}{TOAS07_p2_phy_s3_para3_4_1.png}
\end{bankpassage}
\begin{bankq}[id=TOAS07-P2-PHY-S3-Q03, source={TOAS - Test 07 (01 May 2022)}, paper=2, qno=3, type=numerical, subtopic={Centre of mass}, answer={-46.00}, answer_src=key, difficulty=3, diff_parts={K1 S2 I0 C0 F1}, diff_note={concepts: momentum of a system under forces | steps: net external force; impulse up to 2 s; components}, group=TOAS07-P2-PHY-S3-G3-4]
The momentum of the system at $t=2$ sec in kg m/s is $\left(x\hat i+y\hat j\right)$. Find $(x+y)(g=10$ m/s$^2)$
\end{bankq}
\begin{bankq}[id=TOAS07-P2-PHY-S3-Q04, source={TOAS - Test 07 (01 May 2022)}, paper=2, qno=4, type=numerical, subtopic={Centre of mass}, answer={19.67}, answer_src=key, difficulty=3, diff_parts={K1 S2 I0 C0 F1}, diff_note={concepts: centre of mass motion | steps: position of each body at t = 1 s; centre of mass; sum of components}, group=TOAS07-P2-PHY-S3-G3-4]
The position vector of centre of mass of the system at $t=1$ s is $\left(x\hat i+y\hat j\right)$. Find $(x+y)$.
\end{bankq}
\begin{bankpassage}[id=TOAS07-P2-PHY-S3-G5-6, source={TOAS - Test 07 (01 May 2022)}, questions={5-6}]
Three balls A, B and C of identical size having uniform mass distribution are placed on a smooth horizontal surface. Given that $m_A=2m$, $m_B=m$ and $m_C=m$. Ball A collides with initial velocity $V$ with B as shown in figure. (Assume all collision to be head on). The coefficient of restitution $(e)=\dfrac12$ for all collisions
\bfig{235.6}{TOAS07_p2_phy_s3_para5_6_1.png}
\end{bankpassage}
\begin{bankq}[id=TOAS07-P2-PHY-S3-Q05, source={TOAS - Test 07 (01 May 2022)}, paper=2, qno=5, type=numerical, subtopic={Collisions}, answer={2.00}, answer_src=key, difficulty=2, diff_parts={K1 S1 I0 C0 F1}, diff_note={concepts: elastic collision | steps: conservation of momentum and energy; speed of A after the collision}, group=TOAS07-P2-PHY-S3-G5-6]
The speed of A after collision of A and B for first time is $V/n$. Find $n$.
\end{bankq}
\begin{bankq}[id=TOAS07-P2-PHY-S3-Q06, source={TOAS - Test 07 (01 May 2022)}, paper=2, qno=6, type=numerical, subtopic={Collisions}, answer={4.00}, answer_src=key, difficulty=2, diff_parts={K1 S1 I0 C0 F1}, diff_note={concepts: elastic collision | steps: conservation of momentum and energy; speed of B after the collision with C}, group=TOAS07-P2-PHY-S3-G5-6]
The speed of B after collision of B and C for first time is $V/n$. Find $n$.
\end{bankq}
\begin{bankq}[id=TOAS07-P2-PHY-S4-Q09, source={TOAS - Test 07 (01 May 2022)}, paper=2, qno=9, type=numerical, subtopic={Centre of mass}, answer={3}, answer_src=key, difficulty=3, diff_parts={K1 S2 I0 C0 F1}, diff_note={concepts: spring between two masses | steps: acceleration of the 20 kg mass; spring force from the 10 kg mass; extension; stored energy}]
Two masses 10kg and 20 kg are connected by a massless spring. A force of 200 N acts on 20 kg mass. At the instant when the 10 kg mass has an acceleration 12 m/s$^2$ the energy stored (in J) in the spring ($k=2400$ N/m) is.
\bfig{119.6}{TOAS07_p2_phy_s4_q09_1.png}
\end{bankq}
\begin{bankq}[id=TOAS08-P1-PHY-S1-Q02, source={TOAS - Test 08 (15 May 2022)}, paper=1, qno=2, type=multi, subtopic={Collisions}, answer={A,B,C}, answer_src=key, difficulty=7, diff_parts={K2 S3 I2 C0 F1}, diff_note={concepts: elastic collisions of a block, wedge with a wall | steps: momentum and energy conservation on the first slide; velocity of the wedge after the slide; repeated collisions with the wall; final states for each range of n; compare with the options | I: sequence of collisions ends when velocities stop reversing}]
All the surface are smooth and collisions are elastic in nature. The system is released from rest. Then
\bfig{146.3}{TOAS08_p1_phy_s1_q02_1.png}
\begin{qoptions}
\opt{A}{When block first time slide down wedge and reaches horizontal surface velocity of centre of mass of block-wedge system have zero velocity.}
\opt{B}{After a long time for $n<3$ velocity of centre of mass of block wedge system will be $\sqrt{\dfrac{8ghn}{(n+1)^3}}$.}
\opt{C}{After a long time for $4>n>3$ velocity of centre of mass of block wedge system will be $\sqrt{\dfrac{32ghn(n-1)^2}{(n+1)^5}}$.}
\opt{D}{Velocity of centre of mass of wedge system will not have vertical component.}
\end{qoptions}
\end{bankq}
\begin{bankq}[id=TOAS08-P1-PHY-S3-Q07, source={TOAS - Test 08 (15 May 2022)}, paper=1, qno=7, type=numerical, subtopic={Collisions}, answer={1.50}, answer_src=key, difficulty=6, diff_parts={K2 S3 I1 C0 F1}, diff_note={concepts: collision with restitution, friction impulse | steps: speed at impact; vertical rebound speed from e; friction impulse and horizontal velocity; flight time; horizontal distance | I: friction impulse is limited}]
A ball mass $m=1$kg and radius $R=10$cm is rotating with angular velocity $\omega=20$ rad/s about horizontal axis. It dropped from a height $h=1.1$m on a rough horizontal surface ($\mu_k=0.5$) as shown in figure. If coefficient of resistitution for collision is $e=0.5$ then find horizontal distance travelled by ball between points of $1^{\text{st}}$ and $2^{\text{nd}}$ collision with ground (in m)
\bfig{183.5}{TOAS08_p1_phy_s3_q07_1.png}
\end{bankq}
\begin{bankq}[id=TOAS08-P2-PHY-S2-Q08, source={TOAS - Test 08 (15 May 2022)}, paper=2, qno=8, type=single, subtopic={Collisions}, answer={B}, answer_src=key, difficulty=3, diff_parts={K2 S2 I0 C0 F0}, diff_note={concepts: collision geometry, impulse | steps: geometry of the contact with the radius ratio; impulse components; speed after collision}]
A ball of radius $r=(\sqrt2-1)R$ falls on both balls symmetrically as shown with speed 10 m/s. If $A$ \& $B$ moves with 3 m/s just after collision then match the List--I with list--II. All quantities in List--II are in SI units.
\bfig{91.6}{TOAS08_p2_phy_s2_q08_1.png}
\begin{center}\small\renewcommand{\arraystretch}{1.1}\setlength{\extrarowheight}{3.5pt}\setlength{\tabcolsep}{4pt}
\begin{tabular}{|c|l|c|c|}
\hline
 & \multicolumn{1}{c|}{List-I} & & List-II\\[3pt]
\hline
(P) & \parbox[t]{4.2cm}{Magnitude of impulse of normal between $A$ \& $C$.} & (1) & 6\\[3pt]
(Q) & \parbox[t]{4.2cm}{Magnitude of impulse of tension in any string.} & (2) & 0.2\\[3pt]
(R) & \parbox[t]{4.2cm}{Velocity of $C$ just after collision.} & (3) & $6\sqrt2$\\[3pt]
(S) & \parbox[t]{4.2cm}{Magnitude of net impulse on $C$ during collision.} & (4) & 12\\[3pt]
 & & (5) & 2\\[3pt]
\hline
\end{tabular}
\end{center}
\begin{qoptions}
\opt{A}{$P\to1;Q\to3;R\to1;S\to2$}
\opt{B}{$P\to3;Q\to1;R\to5;S\to4$}
\opt{C}{$P\to1;Q\to1;R\to2;S\to4$}
\opt{D}{$P\to1;Q\to1;R\to2;S\to2$}
\end{qoptions}
\end{bankq}
\begin{bankq}[id=TOAS08-P2-PHY-S2-Q09, source={TOAS - Test 08 (15 May 2022)}, paper=2, qno=9, type=single, subtopic={Variable mass}, answer={B}, answer_src=key, difficulty=2, diff_parts={K1 S2 I0 C0 F0}, diff_note={concepts: carpet being laid and variable mass | steps: mass in motion as a function of time; force from the momentum change; power and work}]
A long thin carpet of length 4m is being laid on the floor. One end of the carpet is bent and then pulled backward with constant velocity $v=2$m/s above the part of carpet which is still at rest on the floor. Carpet has mass density of 0.5 kg/m. All quantities in List--II in SI units.
\bfig{183.7}{TOAS08_p2_phy_s2_q09_1.png}
\begin{center}\small\renewcommand{\arraystretch}{1.1}\setlength{\extrarowheight}{3.5pt}\setlength{\tabcolsep}{4pt}
\begin{tabular}{|c|l|c|c|}
\hline
 & \multicolumn{1}{c|}{List-I} & & List-II\\[3pt]
\hline
(P) & \parbox[t]{4.2cm}{Velocity of com of part of carpet still at rest.} & (1) & 2\\[3pt]
(Q) & \parbox[t]{4.2cm}{Velocity of com of moving part of carpet.} & (2) & 1.5\\[3pt]
(R) & \parbox[t]{4.2cm}{Mass being added in moving part per sec.} & (3) & 1\\[3pt]
(S) & \parbox[t]{4.2cm}{Force needed to pull the moving part.} & (4) & 0.5\\[3pt]
 & & (5) & 0\\[3pt]
\hline
\end{tabular}
\end{center}
\begin{qoptions}
\opt{A}{$P\to5;Q\to2;R\to3;S\to1$}
\opt{B}{$P\to4;Q\to2;R\to4;S\to1$}
\opt{C}{$P\to5;Q\to2;R\to4;S\to3$}
\opt{D}{$P\to3;Q\to3;R\to3;S\to1$}
\end{qoptions}
\end{bankq}
\begin{bankq}[id=TOAS08-P2-PHY-S3-Q05, source={TOAS - Test 08 (15 May 2022)}, paper=2, qno=5, type=numerical, subtopic={Collisions}, answer={0.60}, answer_src=key, difficulty=5, diff_parts={K2 S2 I1 C0 F1}, diff_note={concepts: collision of a particle with a bowl, rigid body | steps: momentum conservation; angular momentum about the centre of mass; velocity of the particle in the vertical direction | I: use angular momentum about the combined centre of mass}]
A hollow hemispherical bowl of mass $m$ is placed on a smooth surface as shown. A particle of same mass strikes it and sticks with the bowl at $A$. Find velocity of particle just after collision in vertical direction ? ($u=5$ m/s)
\bfig{127.7}{TOAS08_p2_phy_s3_q05_1.png}
\end{bankq}
\begin{bankq}[id=TOAS08-P2-PHY-S3-Q06, source={TOAS - Test 08 (15 May 2022)}, paper=2, qno=6, type=numerical, subtopic={Momentum}, answer={0.67}, answer_src=key, difficulty=2, diff_parts={K1 S1 I0 C0 F1}, diff_note={concepts: pressure of an absorbed electron beam | steps: momentum flux P = n m v\^{}2; current = n e v A}]
A beam of fast moving electrons having cross-sectional area $A=1$ cm$^2$ falls normally on a flat surface. The electrons are absorbed by the surface and the average pressure exerted by the electrons on this surface is found to be $P=9.1$ Pa. If the electrons are moving with a speed $v=24\times10^7$ m/s, then find the effective current (in A) through any cross-section of the electron beam. (mass of electron $=9.1\times10^{-31}$ kg)
\end{bankq}
\begin{bankq}[id=TOAS19-P1-PHY-S1-Q01, source={TOAS - Test 19 (05 Feb 2023)}, paper=1, qno=1, type=multi, subtopic={Centre of mass}, answer={B,C,D}, answer_src=key, difficulty=6, diff_parts={K2 S2 I1 C1 F1}, diff_note={concepts: motion of the centre of mass, internal forces | steps: centre of mass of insect and disc stays at rest; positions of disc centre and insect about it; relative angular velocity; force on the disc}, also_in={TNAS - Test 01 (05 Feb 2023) P2 Q1}]
On a smooth horizontal surface there is a disc of mass $2m$ and radius $R$ is placed. An insect of mass $m$ sits gently at a distance $\dfrac R2$ from center of disc and crawls uniformly with speed `$v_0$' with respect to centre of disc along circular path with respect to centre of disc.
\begin{qoptions}
\opt{A}{With respect to a ground observer insect moves in circular path of radius $\dfrac R2$.}
\opt{B}{With respect to a ground observer centre of disc moves on a circular path of radius $\dfrac R6$.}
\opt{C}{Net external force on this system is zero.}
\opt{D}{Insect applies a force $\dfrac4{27}mv_0^2R$ on disc tangentially to the surface of disc.}
\end{qoptions}
\end{bankq}
\begin{bankq}[id=TOAS19-P1-PHY-S3-Q04, source={TOAS - Test 19 (05 Feb 2023)}, paper=1, qno=4, type=numerical, subtopic={Variable mass}, answer={10.00}, answer_src=key, difficulty=7, diff_parts={K2 S2 I2 C1 F1}, diff_note={concepts: chain over a pulley with pick-up, drop | steps: energy lost in picking up links; energy lost on dropping; gravity balance; height}]
The ends of a chain lie in piles at A and C of negligible height. When given an initial speed 10 m/s, the chain keeps on moving at that speed over the fixed pulley B. Neglect friction, Find $h$ in meter. Assume chain do not tangle.
\bfig{58.2}{TOAS19_p1_phy_s3_q04_1.png}
\end{bankq}

\section*{TOPS papers}
\begin{bankq}[id=TOPS02-P1-PHY-S3-Q03, source={TOPS - Test 02 (28 Aug 2022)}, paper=1, qno=3, type=numerical, subtopic={Momentum}, answer={5.00}, answer_src=key, difficulty=7, diff_parts={K2 S2 I2 C1 F1}, diff_note={concepts: momentum of a chain, geometry | steps: length of the chain on the slope and on the flat; velocity direction of each part; vector sum of the momentum | chain around a hill}]
A uniform chain of mass per unit length $\lambda$ is sliding down the hill with constant speed $u$. Dimension of the hill are $h$ and $d$ as shown. Find the momentum (in kg -- m/s) of the chain at the instant shown. (Given $\lambda=1$ kg/m, $d=3$m, $h=4$m, and $u=1$m/s)
\bfig{66.1}{TOPS02_p1_phy_s3_q03_1.png}
\end{bankq}
\begin{bankq}[id=TOPS02-P2-PHY-S1-Q04, source={TOPS - Test 02 (28 Aug 2022)}, paper=2, qno=4, type=multi, subtopic={Momentum}, answer={A,B,D}, answer_src=key, difficulty=10, diff_parts={K3 S3 I2 C2 F1}, diff_note={concepts: momentum conservation, curvature, wedge | steps: conserve horizontal momentum; accelerations of block and wedge; velocity at the top; velocity relative; radius of curvature | moving wedge}]
A wedge of mass $m$ is kept on a smooth table and its inclined surface is also smooth. A small block of mass $m$ is projected from the bottom along the incline surface with velocity $u$. Assume that the block remains on the incline and take $\theta=45^\circ$, $g=10$ m/s$^2$.
\bfig{159.2}{TOPS02_p2_phy_s1_q04_1.png}
\begin{qoptions}
\opt{A}{Magnitude of acceleration of block is more than magnitude of acceleration of wedge.}
\opt{B}{The block strikes the table normally.}
\opt{C}{Radius of curvature of the path of the block when it is at the highest point is $u^2/16g$.}
\opt{D}{Radius of curvature of the path of the block when it is at the highest point is $3u^2/16g$.}
\end{qoptions}
\end{bankq}
\begin{bankq}[id=TOPS02-P2-PHY-S3-Q03, source={TOPS - Test 02 (28 Aug 2022)}, paper=2, qno=3, type=numerical, subtopic={Collisions}, answer={5.00}, answer_src=key, difficulty=9, diff_parts={K3 S3 I2 C1 F1}, diff_note={concepts: restitution, geometric series, kinematics | steps: speeds after each bounce; heights as a geometric series; total distance; total time; average speed; 2v | infinite bounces}]
A rubber ball is released from a height of 5 m above the floor. It bounces back repeatedly, always rising to 0.81 times of the height through which it falls. The average speed of the ball is $v$ (in m/s), then value of $2v$ to nearest integer is : (Take $g=10$ ms$^{-2}$)
\end{bankq}
\begin{bankq}[id=TOPS03-P1-PHY-S1-Q03, source={TOPS - Test 03 (20 Nov 2022)}, paper=1, qno=3, type=single, subtopic={Momentum}, answer={B}, answer_src=key, difficulty=8, diff_parts={K3 S3 I1 C2 F0}, diff_note={concepts: momentum conservation, energy conservation, collision | steps: momentum in the collision; common velocity condition at maximum deflection; horizontal momentum with the trolley; energy equation; angle | shell, ball and trolley}]
A ball of mass $m=1$ kg is hung vertically by a thread of length $\ell=1.50$ m. Upper end of the thread is attached to the ceiling of a trolley of mass $M=4$ kg. Initially trolley is stationary and it is free to move along horizontal rails without friction. A shell of mass $m=1$ kg, moving horizontally with velocity $v_0=6$ ms$^{-1}$, collides with the ball and gets stuck with it. As a result, thread starts to deflect towards right. Calculate its maximum deflection with the vertical. ($g=10$ ms$^{-2}$)
\bfig{143.7}{TOPS03_p1_phy_s1_q03_1.png}
\begin{qoptions}
\opt{A}{$30^\circ$}
\opt{B}{$37^\circ$}
\opt{C}{$53^\circ$}
\opt{D}{$60^\circ$}
\end{qoptions}
\end{bankq}
\begin{bankq}[id=TOPS03-P1-PHY-S4-Q09, source={TOPS - Test 03 (20 Nov 2022)}, paper=1, qno=9, type=numerical, subtopic={Momentum}, answer={6}, answer_src=key, difficulty=10, diff_parts={K3 S3 I2 C2 F1}, diff_note={concepts: momentum conservation, energy conservation, helical motion | steps: horizontal momentum for the cylinder and particle; condition at the top; energy equation; relative speed; speed in terms of Rg | helical groove}]
Figure shows a hollow vertical cylinder of mass `$M$', radius `$R$' and height `$R$', kept on a smooth horizontal floor. It has smooth helical groove on the inner surface. Number of turns in the helical groove is large, so that any jerk can be neglected, when a particle of mass `$M$' enters the groove through a small hole near the bottom in the cylinder moving along the tangent on the floor. Take $R=10$cm and $g=10$m/s$^2$.
\bfig{82.5}{TOPS03_p1_phy_s4_q09_1.png}
The minimum speed $v_0$ of the particle on the horizontal plane, so that it reaches the top of the cylinder is $\sqrt{KRg}$. Find $K$.
\end{bankq}
\begin{bankq}[id=TOPS03-P2-PHY-S2-Q01, source={TOPS - Test 03 (20 Nov 2022)}, paper=2, qno=1, type=numerical, subtopic={Centre of mass}, answer={2.66}, answer_src=key, difficulty=8, diff_parts={K3 S2 I2 C1 F1}, diff_note={concepts: constraint, spring force, Newton's laws | steps: forces on each body; condition of maximum deformation; accelerations equal; solve for the deformation | spring between two blocks}]
In the diagram shown, there is no friction at any contact surface. Initially, the spring has no deformation. What will be the maximum deformation in the spring if a constant horizontal force $F$ is applied as shown? Consider all the strings to be sufficiently large. The spring constant is $K$. If the maximum deformation in the spring is $QF/K$, find $Q$.
\bfig{192.0}{TOPS03_p2_phy_s2_q01_1.png}
\end{bankq}
\begin{bankq}[id=TOPS03-P2-PHY-S2-Q02, source={TOPS - Test 03 (20 Nov 2022)}, paper=2, qno=2, type=numerical, subtopic={Centre of mass}, answer={0.33}, answer_src=key, difficulty=8, diff_parts={K2 S2 I2 C2 F1}, diff_note={concepts: centre of mass, spring forces | steps: centre of mass accelerates as in free fall; extension of each spring at equilibrium; natural length; distance from the second ball | three balls with springs}]
Three small balls of equal mass($m$) are suspended from a thread and two springs of same force constant ($k$) such that the distance between the first and the second ball and the second and third ball are equal. Thus the centre of mass of the whole system coincides with the second ball. The thread supporting the upper ball is cut and system starts a free fall. Then the distance of the centre of mass of the system from the second ball when both the springs acquire their natural length in the falling system is $Nmg/k$. Find $N$.
\bfig{39.7}{TOPS03_p2_phy_s2_q02_1.png}
\end{bankq}
\begin{bankq}[id=TOPS03-P2-PHY-S2-Q03, source={TOPS - Test 03 (20 Nov 2022)}, paper=2, qno=3, type=numerical, subtopic={Centre of mass}, answer={2.00}, answer_src=key, difficulty=8, diff_parts={K3 S2 I2 C1 F1}, diff_note={concepts: centre of mass, friction, relative motion | steps: forces on each block; acceleration of each; centre of mass acceleration | blocks with a string}]
Both the blocks shown in the given arrangement are given together a horizontal velocity towards right. If $x$ ms$^{-2}$ be the subsequent acceleration of centre of mass of the system of blocks, find the value of $x$ (in ms$^{-2}$). Given, $m_A=2m_B$ and $g=10$ ms$^{-2}$, all in SI units.
\bfig{139.8}{TOPS03_p2_phy_s2_q03_1.png}
\end{bankq}
\begin{bankq}[id=TOPS03-P2-PHY-S2-Q04, source={TOPS - Test 03 (20 Nov 2022)}, paper=2, qno=4, type=numerical, subtopic={Centre of mass}, answer={1.81}, answer_src=solved, working={Separation on the dome at $\cos\theta=2/3$. Particle at height $2R/3$, $x=R\sqrt5/3$; hollow hemisphere centre of mass $R/2$ above $O$. Centre of mass of the pair: $y=\dfrac{(2/3+2)R}5=\dfrac{8R}{15}$, $x=\dfrac{\sqrt5R}{15}$, distance $\dfrac{\sqrt{69}R}{15}=\dfrac R\alpha$, $\alpha\approx1.81$}, difficulty=10, diff_parts={K3 S3 I2 C2 F1}, diff_note={concepts: centre of mass, circular motion, condition for separation | steps: condition of separation on the dome; position of the particle; centre of mass of the hollow hemisphere; combined centre of mass; distance | hollow dome}]
A small particle of mass $m$ starts sliding down from rest along the smooth surface of a fixed hollow hemisphere of mass ($M=4m$). The distance of centre of mass of (particle $+$ hemisphere) from centre $O$ of hemisphere, when the particle separates from the surface of hemisphere is ($R/\alpha$). Find the value of $\alpha$.
\bfig{149.6}{TOPS03_p2_phy_s2_q04_1.png}
\end{bankq}
\begin{bankq}[id=TOPS03-P2-PHY-S2-Q05, source={TOPS - Test 03 (20 Nov 2022)}, paper=2, qno=5, type=numerical, subtopic={Momentum}, answer={1.50}, answer_src=key, difficulty=10, diff_parts={K3 S3 I2 C2 F1}, diff_note={concepts: momentum conservation, energy conservation, wedge | steps: wedge held by the wall; momentum condition; energy; common velocity at the highest point; height | wedge against a wall}]
For the figure shown, the wedge $+$ particle is resting at a vertical wall. The particle of mass $\left(m=\dfrac M2\right)$ is released from rest from the position shown. The wedge has mass $M$. The maximum height reached by the particle on the circular groove is ($R/\alpha$). Find the value of $\alpha$.
\bfig{141.3}{TOPS03_p2_phy_s2_q05_1.png}
\end{bankq}
\begin{bankq}[id=TOPS03-P2-PHY-S3-Q03, source={TOPS - Test 03 (20 Nov 2022)}, paper=2, qno=3, type=numerical, subtopic={Momentum}, answer={3}, answer_src=key, difficulty=9, diff_parts={K2 S3 I2 C2 F1}, diff_note={concepts: elastic rope, momentum and relative motion | steps: rope taut condition; velocity of the centre of mass; circular motion in the frame; angle turned in the time; displacement | rope becomes taut at an angle}]
Two identical particles $A$ and $B$ connected by an inextensible, ideal, massless, elastic rope of length $2\ell$ are placed at rest on a smooth horizontal floor at a separation of $\sqrt3\ell$. At $t=0$, $A$ is given a horizontal velocity $u$ perpendicular to the line joining $AB$ as shown in the figure. If the distance moved by $A$ till the time $t=\dfrac{11\ell}u$ is $\sqrt3\left(2+\sqrt x\right)\ell$, where $x$ is a positive integer, find the value of $x$.
\bfig{36.2}{TOPS03_p2_phy_s3_q03_1.png}
\end{bankq}
\begin{bankq}[id=TOPS03-P2-PHY-S3-Q04, source={TOPS - Test 03 (20 Nov 2022)}, paper=2, qno=4, type=numerical, subtopic={Momentum}, answer={1}, answer_src=key, difficulty=10, diff_parts={K3 S3 I2 C2 F1}, diff_note={concepts: momentum conservation, energy conservation, circular motion | steps: horizontal momentum; energy equation; normal force zero condition; relative speed; polynomial in cos theta; coefficients | particle leaves a moving hemisphere}]
A hemisphere of mass $2m$ is free to slide with its base on a smooth horizontal table. A particle of mass $m$ is placed on the hemisphere at an angular distance $\alpha=\dfrac\pi3$ from the top of the hemisphere. The radius to the point of contact at which the particle leaves the surface makes an angle $\theta$ with the vertical given by the equation $\cos3\theta+3a\cos\theta+b=0$, where $a$ and $b$ are integers. Find $a+b$.
\end{bankq}
\begin{bankq}[id=TOPS04-P1-PHY-S1-Q03, source={TOPS - Test 04 (18 Dec 2022)}, paper=1, qno=3, type=multi, subtopic={Collisions}, answer={A,B,D}, answer_src=key, difficulty=10, diff_parts={K3 S3 I2 C2 F1}, diff_note={concepts: collision, circular motion, energy | steps: speed after the first collision; tension; height reached; second collision and horizontal circle; period}]
A pendulum bob has mass $m$. The length of pendulum is $\ell$. It is initially at rest. A particle $P$ of mass $m/2$ moving horizontally along $-$ve $x$-direction with velocity $\sqrt{2g\ell}$ collides with the bob and comes to rest. When the bob comes to rest momentarily, another particle $Q$ of mass $m$ moving horizontally along $z$-direction collides with the bob and sticks to it. It is observed that the bob now moves along a horizontal circle.

The floor is a horizontal surface at a distance $2\ell$ below the point of suspension of the pendulum.
\bfig{205.5}{TOPS04_p1_phy_s1_q03_1.png}
\begin{qoptions}
\opt{A}{Tension in string immediately after the first collision is $\dfrac32mg$.}
\opt{B}{The height of circular path above the floor is $\dfrac{5\ell}4$}
\opt{C}{Time period of the circular motion $2\pi\sqrt{\dfrac\ell g}$}
\opt{D}{Time period of the circular motion $2\pi\sqrt{\dfrac{3\ell}{4g}}$}
\end{qoptions}
\end{bankq}
\begin{bankpassage}[id=TOPS04-P1-PHY-S2-G9-10, source={TOPS - Test 04 (18 Dec 2022)}, questions={9-10}]
A light rod of length $L$ is hinged to a plank of mass $m$. The plank is lying on the edge of a horizontal table such that the rod can swing freely in the vertical plane without any hindrance from the table. A particle of mass $m$ is attached to the end of the rod and system is released from $\theta=0^\circ$ position (see figure)
\bfig{107.2}{TOPS04_p1_phy_s2_para9_10_1.png}
\end{bankpassage}
\begin{bankq}[id=TOPS04-P1-PHY-S2-Q09, source={TOPS - Test 04 (18 Dec 2022)}, paper=1, qno=9, type=single, subtopic={Momentum}, answer={D}, answer_src=key, difficulty=6, diff_parts={K2 S2 I2 C1 F0}, diff_note={concepts: momentum conservation, energy | steps: horizontal momentum for plank and particle; energy conservation; speed of the plank at 180 degrees}, group=TOPS04-P1-PHY-S2-G9-10]
Assume that friction is absent everywhere and calculate the speed of the plank when the rod makes $\theta=180^\circ$.
\begin{qoptions}
\opt{A}{$\sqrt{8gL}$}
\opt{B}{$\sqrt{6gL}$}
\opt{C}{$2\sqrt{gL}$}
\opt{D}{$\sqrt{2gL}$}
\end{qoptions}
\end{bankq}
\begin{bankq}[id=TOPS04-P1-PHY-S2-Q10, source={TOPS - Test 04 (18 Dec 2022)}, paper=1, qno=10, type=single, subtopic={Momentum}, answer={A}, answer_src=key, difficulty=9, diff_parts={K3 S3 I2 C2 F0}, diff_note={concepts: torque, energy, normal force | steps: energy and speed at angle; forces on the plank; normal force as a function of angle; minimise; value}, group=TOPS04-P1-PHY-S2-G9-10]
Assume that friction between the plank and the table is large enough to prevent it from slipping and calculate the smallest normal force applied by the plank on the table.
\begin{qoptions}
\opt{A}{$\dfrac{2mg}3$}
\opt{B}{$\dfrac{mg}3$}
\opt{C}{$mg$}
\opt{D}{$\dfrac{mg}2$}
\end{qoptions}
\end{bankq}
\begin{bankq}[id=TOPS04-P2-PHY-S1-Q02, source={TOPS - Test 04 (18 Dec 2022)}, paper=2, qno=2, type=single, subtopic={Momentum}, answer={B}, answer_src=key, difficulty=5, diff_parts={K2 S2 I1 C1 F0}, diff_note={concepts: impulse, average force | steps: velocity change; time on the arc; average force}]
A fixed U-shaped smooth wire has a semi-circular bending between $A$ and $B$ as shown in the figure. A bead of mass `$m$' moving with uniform speed $v$ through the wire enters the semicircular bend at $A$ and leaves at $B$. The magnitude of average force exerted by the bead on the part $AB$ of the wire is :-
\bfig{108.3}{TOPS04_p2_phy_s1_q02_1.png}
\begin{qoptions}
\opt{A}{0}
\opt{B}{$\dfrac{4mv^2}{\pi d}$}
\opt{C}{$\dfrac{2mv^2}{\pi d}$}
\opt{D}{None of these}
\end{qoptions}
\end{bankq}
\begin{bankq}[id=TOPS04-P2-PHY-S1-Q03, source={TOPS - Test 04 (18 Dec 2022)}, paper=2, qno=3, type=single, subtopic={Centre of mass}, answer={C}, answer_src=key, difficulty=6, diff_parts={K2 S2 I2 C1 F0}, diff_note={concepts: centre of mass, infinite series | steps: mass ratio and sizes; positions of each centre; infinite sum; coordinate}]
An infinite number of bricks are placed one over the other as shown in the figure. Each succeeding brick having half the length and breadth of its preceding brick and the mass of each succeeding bricks being $1/4^{\text{th}}$ of the preceding one. Take `$O$' as the origin, the $x$-coordinate of centre of mass of the system of bricks is at :-
\bfig{154.5}{TOPS04_p2_phy_s1_q03_1.png}
\begin{qoptions}
\opt{A}{$-\dfrac a7$}
\opt{B}{$\dfrac{3a}7$}
\opt{C}{$-\dfrac{3a}7$}
\opt{D}{$-\dfrac{2a}7$}
\end{qoptions}
\end{bankq}
\begin{bankq}[id=TOPS04-P2-PHY-S2-Q02, source={TOPS - Test 04 (18 Dec 2022)}, paper=2, qno=2, type=numerical, subtopic={Variable mass}, answer={98}, answer_src=solved, working={Thrust $u\,\dfrac{dm}{dt}\ge Mg$, so $\dfrac{dm}{dt}=\dfrac{3000\times9.8}{300}=98$ kg/s}, difficulty=3, diff_parts={K1 S1 I0 C1 F1}, diff_note={concepts: rocket thrust | steps: thrust equals weight; rate}]
Find the minimum rate (kg/sec) at which the fuel must be consumed by a rocket so as to be able to take off vertically. Given that the total mass at the time of take off is 3000 kg and the velocity of fuel expelled is 300 m/s. [$g=9.8$ m/s$^2$]
\end{bankq}
\begin{bankq}[id=TOPS05-P1-PHY-S1-Q01, source={TOPS - Test 05 (14 Jan 2023)}, paper=1, qno=1, type=multi, subtopic={Momentum}, answer={B,D}, answer_src=key, difficulty=6, diff_parts={K2 S2 I1 C1 F1}, diff_note={concepts: momentum conservation, work-energy | steps: momentum of P from conservation; kinetic energies at the start and at T; work on each particle; sum}]
Two particles, $P$ of mass $2m$ and $Q$ of mass $m$, are subjected to mutual force of attraction and no other force acts on them. At $t=0$, $P$ is at rest at point $O$ and $Q$ is moving away from $O$ with a speed $5u$. At a later instant $t=T$ (before any collision has taken place), $Q$ is moving towards $O$ with speed $u$. Then
\begin{qoptions}
\opt{A}{momentum of particle $P$ at $t=T$ is zero.}
\opt{B}{momentum of particle $P$ at $t=T$ is $6mu$.}
\opt{C}{work done by the force of attraction during $0\le t\le T$ is $12mu^2$.}
\opt{D}{work done by the force of attraction during $0\le t\le T$ is $-3mu^2$.}
\end{qoptions}
\end{bankq}
\begin{bankq}[id=TOPS05-P1-PHY-S3-Q07, source={TOPS - Test 05 (14 Jan 2023)}, paper=1, qno=7, type=numerical, subtopic={Collisions}, answer={0.88}, answer_src=key, difficulty=9, diff_parts={K2 S3 I2 C2 F1}, diff_note={concepts: momentum conservation, ring and beads | steps: conserve momentum; energy conservation; relative velocity at the impact; solve the equations; value}]
Two small beads each of mass $m$ are attached to a smooth ring of mass $m$ as shown in the figure. The whole system is kept on smooth horizontal table. Initially, each bead has given velocity $v_0$ as shown in figure. If the velocity of bead just before impact is $xv_0$, the value of $x$ is.
\bfig{103.0}{TOPS05_p1_phy_s3_q07_1.png}
\end{bankq}
\begin{bankq}[id=TOPS05-P2-PHY-S3-Q02, source={TOPS - Test 05 (14 Jan 2023)}, paper=2, qno=2, type=numerical, subtopic={Variable mass}, answer={5}, answer_src=key, difficulty=9, diff_parts={K2 S3 I2 C2 F1}, diff_note={concepts: variable mass, differential equation | steps: momentum equation with falling sand; equation of motion; integrate; time 200 s; evaluate}]
A container of mass 200 kg is moving towards right by applying a constant force 100 N. Initial velocity of container is zero. Sand is falling in container with constant rate of 1 kg/s with the help of a fixed cone as shown. Velocity of container at $t=200$ sec is $10N$ m/s. Find $N$. ($\ell$n 1.5 = 0.4)
\bfig{112.9}{TOPS05_p2_phy_s3_q02_1.png}
\end{bankq}
\begin{bankq}[id=TOPS05-P2-PHY-S3-Q04, source={TOPS - Test 05 (14 Jan 2023)}, paper=2, qno=4, type=numerical, subtopic={Collisions}, answer={6}, answer_src=key, difficulty=5, diff_parts={K1 S2 I1 C1 F1}, diff_note={concepts: restitution | steps: velocities along the line; relative speed ratio; e}]
Figure shows diagram of two bodies just before and just after collision. If $e$ is the coefficient of restitution, find $10\,e$.
\bfig{100.4}{TOPS05_p2_phy_s3_q04_1.png}
\end{bankq}

\section*{TMAS papers}
\begin{bankq}[id=TMAS02-P1-PHY-S1-Q04, source={TMAS - Test 02 (06 Aug 2023)}, paper=1, qno=4, type=single, subtopic={Collisions}, answer={C}, answer_src=key, difficulty=4, diff_parts={K1 S2 I1 C1 F0}, diff_note={concepts: momentum and energy in a string jerk | steps: common speed along the string; kinetic energy loss; cases}]
Two identical masses $m$ each are connected by a loose string as shown. Masses $A$ and $B$ are moving with velocity $2V$ and $V$ on a smooth horizontal surface. After sufficiently long time.
\bfig{124.2}{TMAS02_p1_phy_s1_q04_1.png}
\begin{qoptions}
\opt{A}{system of masses would have lost a kinetic energy of magnitude $\dfrac{mv^2}{4}$ only if collision between $A$ and $B$ is perfectly inelastic}
\opt{B}{system of masses would not loose any kinetic energy if collision between $A$ and $B$ is elastic}
\opt{C}{system of masses must have lost a kinetic energy $\dfrac{mv^2}{4}$ irrespective of nature of collision between $A$ and $B$.}
\opt{D}{system of masses must have lost all its kinetic energy.}
\end{qoptions}
\end{bankq}
\begin{bankq}[id=TMAS02-P1-PHY-S2-Q08, source={TMAS - Test 02 (06 Aug 2023)}, paper=1, qno=8, type=multi, subtopic={Collisions}, answer={A}, answer_src=key, difficulty=8, diff_parts={K2 S3 I2 C1 F1}, diff_note={concepts: impulsive tension, collision | steps: speed after the collision; string tension impulse; speed of body 3; acceleration; check}]
Consider the system of bodies in figure, their masses are specified in the figure, as well as the length of the string connecting bodies 2 and 3, which are at rest at the same height. Body 1 which is released from rest at a height $h$ makes an inelastic collision with the balance pan. Both strings are considered to be inelastic. The gravitational acceleration $g$ is known (the mass of the balance pan will be neglected)
\bfig{123.7}{TMAS02_p1_phy_s2_q08_1.png}
\begin{qoptions}
\opt{A}{Just after sticking to pan both bodies (1) and (2) will have speed given by $\dfrac{m}{M+m}\sqrt{2gh}$}
\opt{B}{Just after sticking to pan both bodies (1) and (2) will have speed given by $\dfrac{m}{2M+m}\sqrt{2gh}$}
\opt{C}{The acceleration of body 3, just after body 1 makes an inelastic collision with the balance pan is $4g\dfrac{m^2}{(M+m)^2}\dfrac{h}{\ell}$}
\opt{D}{The acceleration of body 3, just after body 1 makes an inelastic collision with the balance pan is $2g\dfrac{m^2}{(2M+m)^2}\dfrac{h}{\ell}$}
\end{qoptions}
\end{bankq}
\begin{bankq}[id=TMAS02-P1-PHY-S2-Q09, source={TMAS - Test 02 (06 Aug 2023)}, paper=1, qno=9, type=multi, subtopic={Variable mass}, answer={A,B}, answer_src=key, difficulty=7, diff_parts={K2 S2 I2 C1 F1}, diff_note={concepts: variable mass, chain collected | steps: mass collected; momentum equation; force; work}]
A massless bucket is initially at rest next to one end of a chain that lies in a straight line on the floor, as shown in Fig. The chain has uniform mass density $\lambda$ (kg/m). You push on the bucket (so that it gathers up the chain) with the force $F(t)$ that gives the bucket and whatever chain is inside, a constant acceleration $a$ at all times. `$t$' is time. Friction is absent
\bfig{145.2}{TMAS02_p1_phy_s2_q09_1.png}
\begin{qoptions}
\opt{A}{$F(t)$ at time $t$ is $\dfrac32\lambda a^2t^2$}
\opt{B}{Work done by $F(t)$ upto time $t$ is $\dfrac38\lambda a^3t^4$}
\opt{C}{$F(t)$ at time $t$ is $\dfrac38\lambda a^2t^2$}
\opt{D}{None}
\end{qoptions}
\end{bankq}
\begin{bankq}[id=TMAS02-P1-PHY-S4-Q07, source={TMAS - Test 02 (06 Aug 2023)}, paper=1, qno=7, type=numerical, subtopic={Variable mass}, answer={4}, answer_src=key, difficulty=6, diff_parts={K1 S2 I2 C1 F1}, diff_note={concepts: variable mass on a conveyor | steps: sand momentum; force from the belt; power or torque; value}]
Sand is transported to the construction site using a conveyor belt. The length of the belt is $\ell$, the angle with respect to the horizontal is $\alpha$ ; the belt is driven by the lower pulley with radius $R$, powered externally. The sand is put onto the belt at a constant rate $\mu$ (kg/s). What is the minimal required torque (in N-m) needed to transport the sand ? The coefficient of friction is larger enough for the sand grains to stop moving immediately after hitting the belt ; take the initial velocity of the sand grains to be zero. (Take $\mu=0.2$ kg/sec, $R=1$m, $\ell=20$ m and $\alpha=30^\circ$)
\bfig{130.3}{TMAS02_p1_phy_s4_q07_1.png}
\end{bankq}
\begin{bankq}[id=TMAS02-P2-PHY-S1-Q06, source={TMAS - Test 02 (06 Aug 2023)}, paper=2, qno=6, type=multi, subtopic={Collisions}, answer={A,C}, answer_src=solved, working={Momentum and angular momentum about the centre of mass give $\omega=v_0/(4\ell)$ and kinetic energy after $=\tfrac3{16}mv_0^2$, so energy lost $=\tfrac5{16}mv_0^2$}, difficulty=6, diff_parts={K1 S2 I2 C1 F1}, diff_note={concepts: collision with a dumbbell | steps: momentum components; angular momentum; angular velocity; energy}]
A putty ball of mass $m$ and speed $v_0$ strikes a dumbell as shown in figure. Dumbbell consists of light rod and two particles. Putty sticks to particle after collision. Mark the CORRECT option (s) :-
\bfig{66.3}{TMAS02_p2_phy_s1_q06_1.png}
\begin{qoptions}
\opt{A}{Angular velocity of dumbell is $\dfrac{v_0}{4\ell}$ clockwise.}
\opt{B}{Angular velocity of dumbell is $\dfrac{3v_0}{\ell}$ clockwise.}
\opt{C}{Energy lost in collision is $\dfrac{5mv_0^2}{16}$}
\opt{D}{Energy lost in collision is $\dfrac{3mv_0^2}{16}$}
\end{qoptions}
\end{bankq}
\begin{bankpassage}[id=TMAS02-P2-PHY-S2-G7-8, source={TMAS - Test 02 (06 Aug 2023)}, questions={7-8}]
A system consists of two small spheres of mass $1$kg and $2$kg interconnected by a weightless spring. At $t=0$, the spheres are set in motion with initial velocities as shown in figure. Position vector of centre of mass at $t=0$ is $20\,\hat j$. Initially spring is in its natural length. System moves in earth's gravitational field.
\bfig{108.0}{TMAS02_p2_phy_s2_para7_8_1.png}
\end{bankpassage}
\begin{bankq}[id=TMAS02-P2-PHY-S2-Q07, source={TMAS - Test 02 (06 Aug 2023)}, paper=2, qno=7, type=single, subtopic={Centre of mass}, answer={A}, answer_src=solved, working={$\vec p_0=(18,-4)$ and gravity changes $p_y$ by $-30\times2$, so $\vec p=(18,-64)$ and $x+y=-46$}, difficulty=4, diff_parts={K1 S2 I1 C1 F0}, diff_note={concepts: momentum under gravity | steps: initial momentum; impulse of gravity; value}, group=TMAS02-P2-PHY-S2-G7-8]
The momentum of the system at $t=2$ sec in kg m/s is $\left(x\hat i+y\hat j\right)$. Find $(x+y)$ ($g=10\ \mathrm{m/s^2}$)
\begin{qoptions}
\opt{A}{$-46$}
\opt{B}{$-92$}
\opt{C}{$80$}
\opt{D}{None of these}
\end{qoptions}
\end{bankq}
\begin{bankq}[id=TMAS02-P2-PHY-S2-Q08, source={TMAS - Test 02 (06 Aug 2023)}, paper=2, qno=8, type=single, subtopic={Centre of mass}, answer={A}, answer_src=solved, working={$\vec v_{cm}=(6,-4/3)$; $\vec r=(6,\ 20-\tfrac43-5)$ so $x+y\approx19.67$}, difficulty=3, diff_parts={K1 S1 I1 C1 F0}, diff_note={concepts: motion of the centre of mass | steps: centre of mass velocity; position at 1 s}, group=TMAS02-P2-PHY-S2-G7-8]
The position vector of centre of mass of the system at $t=1$ s is $\left(x\hat i+y\hat j\right)$. Find $(x+y)$.
\begin{qoptions}
\opt{A}{$19.67$}
\opt{B}{$40.3$}
\opt{C}{$67.8$}
\opt{D}{None}
\end{qoptions}
\end{bankq}

\section*{TNAS papers}
\begin{bankq}[id=TNAS01-P2-PHY-S3-Q03, source={TNAS - Test 01 (05 Feb 2023)}, paper=2, qno=3, type=numerical, subtopic={Momentum}, answer={0.665}, answer_src=key, difficulty=7, diff_parts={K1 S2 I2 C2 F1}, diff_note={concepts: recoil on a balanced board | steps: velocity of the second object; times to reach the pivot; condition to tilt}]
There are two small objects on an $L=3$m long wooden board, supported at its middle by a wedge. The masses of the objects are $m_1=0.2$ kg and $m_2=0.3$ kg. Between the objects there is a compressed small spring of negligible mass. The spring is compressed to a length of $\ell=0.3$ m, and tied with a thread. The wedge is exactly below the center of mass of the system. The thread is burnt, and the spring launches the object with mass $m_1$ at a speed of $v_1=2$ m/s. The time (in sec) after which the board will tilt
\bfig{174.2}{TNAS01_p2_phy_s3_q03_1.png}
\end{bankq}
\begin{bankq}[id=TNAS02-P1-PHY-S4-Q07, source={TNAS - Test 02 (07 May 2023)}, paper=1, qno=7, type=numerical, subtopic={Collisions}, answer={6}, answer_src=key, difficulty=7, diff_parts={K2 S2 I2 C1 F1}, diff_note={concepts: elastic collision with a wedge, projectile | steps: rebound direction; velocity components; time of flight}]
A small elastic ball moving horizontally with 40 m/s collides elastically with the smooth stationary wedge of finite mass at point P which is placed on smooth horizontal surface as shown. After that it again strikes on the wedge at the same point P after $t_0$ seconds. Find the value of $t_0$ (in sec). (Take $g=10$ ms$^{-2}$)
\bfig{113.1}{TNAS02_p1_phy_s4_q07_1.png}
\end{bankq}
\begin{bankq}[id=TNAS02-P1-PHY-S4-Q09, source={TNAS - Test 02 (07 May 2023)}, paper=1, qno=9, type=numerical, subtopic={Collisions}, answer={75}, answer_src=key, difficulty=7, diff_parts={K1 S3 I2 C1 F1}, diff_note={concepts: collision of a pendulum bob with a sphere | steps: geometry of contact; impulse direction; restitution; velocities; e}]
A pendulum of mass $m$ connected with string of length $\ell=20\sqrt3$ cm. It is suspended above the centre of a sphere of radius 20 cm and mass $m$, at a distance 40 cm above centre of sphere. After collision pendulum stops. All surfaces are smooth. If coefficient of restitution is $e$. Write value of $100e$.
\bfig{89.6}{TNAS02_p1_phy_s4_q09_1.png}
\end{bankq}
\begin{bankpassage}[id=TNAS02-P2-PHY-S2-G9-10, source={TNAS - Test 02 (07 May 2023)}, questions={9-10}]
A system consists of two small spheres of mass 1kg and 2 kg interconnected by a weightless spring. At $t=0$, the spheres are set in motion with initial velocities as shown in figure. Position vector of centre of mass at $t=0$ is $20\hat j$. Initially spring is in its natural length. System moves in earth's gravitational field.
\bfig{116.5}{TNAS02_p2_phy_s2_para9_10_1.png}
\end{bankpassage}
\begin{bankq}[id=TNAS02-P2-PHY-S2-Q09, source={TNAS - Test 02 (07 May 2023)}, paper=2, qno=9, type=single, subtopic={Centre of mass}, answer={C}, answer_src=key, difficulty=4, diff_parts={K1 S2 I1 C1 F0}, diff_note={concepts: momentum of two masses under gravity | steps: initial momentum; impulse; value}, group=TNAS02-P2-PHY-S2-G9-10]
The momentum of the system at $t=2$ sec in kg m/s is ($g=10$ m/s$^2$)
\begin{qoptions}
\opt{A}{$\left(12\hat i-32\hat j\right)$}
\opt{B}{$\left(16\hat i-64\hat j\right)$}
\opt{C}{$\left(18\hat i-64\hat j\right)$}
\opt{D}{$\left(12\hat i-18\hat j\right)$}
\end{qoptions}
\end{bankq}
\begin{bankq}[id=TNAS02-P2-PHY-S2-Q10, source={TNAS - Test 02 (07 May 2023)}, paper=2, qno=10, type=single, subtopic={Centre of mass}, answer={D}, answer_src=key, difficulty=3, diff_parts={K1 S1 I1 C1 F0}, diff_note={concepts: centre of mass motion | steps: velocity; position at 1 s}, group=TNAS02-P2-PHY-S2-G9-10]
The position vector of centre of mass of the system at $t=1$s is :-
\begin{qoptions}
\opt{A}{$4\hat i+\dfrac{24}{25}\hat j$}
\opt{B}{$6\hat i+\dfrac{39}{2}\hat j$}
\opt{C}{$4\hat i+\dfrac{29}{3}\hat j$}
\opt{D}{$6\hat i+\dfrac{41}{3}\hat j$}
\end{qoptions}
\end{bankq}
\begin{bankq}[id=TNAS02-P2-PHY-S4-Q07, source={TNAS - Test 02 (07 May 2023)}, paper=2, qno=7, type=numerical, subtopic={Variable mass}, answer={10}, answer_src=key, difficulty=6, diff_parts={K1 S2 I2 C1 F1}, diff_note={concepts: friction with falling sand | steps: rate of added mass; force of friction; power}]
A trolley is moving with constant velocity of 1 m/s on a straight track due to some external force of 20N. Sand is continuously falling on it with negligible velocity and moves with trolley. What is magnitude of rate of work done by friction between sand and trolley?
\end{bankq}
\begin{bankq}[id=TNAS02-P2-PHY-S4-Q08, source={TNAS - Test 02 (07 May 2023)}, paper=2, qno=8, type=numerical, subtopic={Momentum}, answer={4}, answer_src=key, difficulty=6, diff_parts={K1 S2 I2 C1 F1}, diff_note={concepts: block with a semi cylindrical hollow | steps: momentum conservation; energy conservation; normal force at the bottom; solve M over m}]
A semi-cylinder is cut out of a block and the block is kept on the horizontal surface. A small body of mass $m$ is released at the top of the cylinder shaped hole. (Friction is negligible).\\
If the block is able to move on the horizontal surface without friction then the force exerted on the block by the small body at the lowest point is $7$ mg/2. Find the ratio of $M/m$ ?
\bfig{89.7}{TNAS02_p2_phy_s4_q08_1.png}
\end{bankq}

\section*{Star Review (Other) papers}
\begin{bankq}[id=STAR01-P1-PHY-S1-Q01, source={Star Review (Other) - Test 01 (12 Sep 2022)}, paper=1, qno=1, type=multi, subtopic={Centre of mass}, answer={B,D}, answer_src=key, difficulty=5, diff_parts={K1 S2 I1 C1 F1}, diff_note={concepts: centre of mass of joined triangular plates | steps: centroid of one triangle; position of the second; area weighted average; both coordinates}]
Two identical uniform equilateral triangular plates are joined as shown. If $(x,y)$ are the coordinates of centre of mass of this combined plate, then :
\bfig{121.7}{STAR01_p1_phy_s1_q01_1.png}
\begin{qoptions}
\opt{A}{$x=\dfrac{\ell\sqrt3}2$}
\opt{B}{$x=\dfrac\ell2$}
\opt{C}{$y=\dfrac\ell2$}
\opt{D}{$y=\dfrac{\ell\sqrt3}2$}
\end{qoptions}
\end{bankq}
\begin{bankq}[id=STAR01-P2-PHY-S1-Q12, source={Star Review (Other) - Test 01 (12 Sep 2022)}, paper=2, qno=12, type=single, subtopic={Centre of mass}, answer={D}, answer_src=key, difficulty=1, diff_parts={K1 S1 I0 C0 F0}, diff_note={concepts: position of the centre of mass | steps: definition and examples}]
The centre of mass of a body :
\begin{qoptions}
\opt{A}{Lies at geometric centre}
\opt{B}{Lies always inside the body.}
\opt{C}{Lies always outside the body.}
\opt{D}{Lies within or outside the body.}
\end{qoptions}
\end{bankq}
\end{document}
